\documentclass{article}

\usepackage{booktabs}
\usepackage{tabularx}
\usepackage{hyperref}
\usepackage{longtable}
\usepackage{pdflscape}
\usepackage{array}

\hypersetup{
    colorlinks=true,
    linkcolor=red,
    citecolor=green,
    filecolor=magenta,
    urlcolor=cyan
}

\title{Hazard Analysis\\\progname}

\author{\authname}

\date{}

\input{../Comments.text}
\input{../Common.text}

\begin{document}

\maketitle
\thispagestyle{empty}

~\newpage

\pagenumbering{roman}

\begin{table}[hp]
\caption{Revision History} \label{TblRevisionHistory}
\begin{tabularx}{\textwidth}{llX}
\toprule
\textbf{Date} & \textbf{Developer(s)} & \textbf{Change}\\
\midrule
Oct 14, 2026 & Qusay Qadir & Initial draft and hazard definition\\
Oct 18, 2026 & Tyler Yue & Drafted scope, components, assumptions, FMEA,
safety requirements, and roadmap based on Volere SRS\\
\bottomrule
\end{tabularx}
\end{table}

~\newpage

\tableofcontents

~\newpage

\pagenumbering{arabic}

\section{Introduction}

This document presents the hazard analysis for \progname{}, an IDE-integrated
tool that captures engineering intent from AI-assisted development, maintains a
versioned intent model with traceability to code and tests, and verifies
AI-proposed changes against that model before they take effect.

The following definitions are adapted from the STPA Handbook
(Leveson and Thomas, 2018):

\begin{itemize}
  \item \textbf{Loss:} Something of value to stakeholders that is unacceptable
  to lose. Losses may include loss of mission, loss of reputation, loss or leak
  of sensitive information, or other stakeholder-unacceptable outcomes.
  \item \textbf{Hazard:} A system state or set of conditions that, together with
  a particular set of worst-case environmental conditions, will lead to a loss.
  \item \textbf{System:} A set of components that act together to achieve a
  common goal. A system may contain subsystems and may be part of a larger
  system.
\end{itemize}

In this software setting, a hazard is a state or condition of \progname{} (or of
its interaction with the IDE, AI coding assistant, network, or workspace) that
can lead to failure, unintended behaviour, loss of engineering control, or
other unacceptable outcomes. The analysis identifies such hazards early and
records strategies to eliminate or mitigate them.

\section{Scope and Purpose of Hazard Analysis}

The purpose of this hazard analysis is to identify hazards that may arise when
using \progname{}, analyze their causes and effects, and define mitigations.
Completing the analysis early produces safety and security requirements that
can be reflected in the SRS and design before implementation locks in unsafe
behaviour.

The losses of concern for \progname{} include:
\begin{itemize}
  \item Loss of engineering intent (requirements, constraints, and rationale
  become incomplete, incorrect, or unrecoverable).
  \item Loss of human control over AI-generated changes (code is produced or
  merged without explicit developer approval).
  \item Loss of traceability between intent and implementation (links to code
  and tests become missing, stale, or unjustified).
  \item Incorrect acceptance decisions (the system reports that a change
  satisfies the intent model when it does not, or vice versa).
  \item Leakage of private project content (source code, requirements, or
  design rationale exposed through external LLM or logging services).
  \item Loss of developer trust or wasted effort (false confidence in the
  intent model, or repeated failed verification cycles).
\end{itemize}

Physical harm is not in scope for this tool. Security, data integrity, AI
over-trust, and software-quality harms are in scope.

\section{System Boundaries and Components}

\progname{} runs locally as an IDE extension coordinated with an AI coding
assistant, with optional calls to external LLM services and Git-backed
versioning of the intent model. The system is divided into the following
major components for hazard brainstorming (not a full design):

\begin{enumerate}
  \item \textbf{Intent Capture / Requirement Extraction} --- accepts
  natural-language requests and extracts structured draft requirements,
  quality constraints, and design constraints (SRS FR1--FR3).
  \item \textbf{Requirement Review and Approval} --- presents proposals for
  approve, edit, or reject before code generation (SRS FR4--FR6, G5).
  \item \textbf{Code Generation Control} --- gates the AI coding assistant so
  implementation proceeds only after an approved decision (SRS FR7--FR8).
  \item \textbf{Intent Model Store} --- persists the versioned structured
  intent model (requirements, constraints, rationale) (SRS FR9, FR11, G2).
  \item \textbf{Traceability} --- maintains justified links from intent
  elements to code and tests (SRS FR9, G3).
  \item \textbf{Change Verification} --- checks AI-proposed or generated
  changes against the intent model and reports results (SRS FR10, FR13, G4).
  \item \textbf{IDE Extension / User Interface} --- VS Code / Cursor (and
  related) surfaces for review, status, and developer interaction.
  \item \textbf{External LLM / Network Services} --- remote models or tracing
  services used for extraction or evaluation.
  \item \textbf{Version Control Integration} --- Git / GitHub versioning of the
  intent model alongside the project.
\end{enumerate}

\section{Critical Assumptions}

The following assumptions bound the analysis. Assumptions that would hide
important failure modes are avoided where possible.

\begin{itemize}
  \item The developer using \progname{} is authorized to work on the local
  project and is not intentionally sabotaging the intent model.
  \item The host IDE and operating system provide a writable project workspace.
  \item An AI coding assistant compatible with the product's integration
  surface is available when gated generation is used.
  \item Network connectivity may fail; hazards from offline or failed LLM
  calls are retained in the FMEA rather than assumed away.
  \item Git is available when intent-model versioning features are used; local
  persistence hazards are still analyzed if Git is unavailable.
\end{itemize}

\section{Failure Mode and Effect Analysis}

Table~\ref{TblFMEA} is the Failure Mode and Effect Analysis (FMEA) for
\progname{}. Recommended actions map to safety and security requirements in
Section~\ref{sec:SR} and/or to existing SRS functional requirements. Hazard
identifiers (for example H1-1) support traceability.

\begin{landscape}
\begin{small}
\begin{longtable}{|p{2.1cm}|p{2.4cm}|p{3.2cm}|p{3.2cm}|p{4.0cm}|p{1.6cm}|}
\caption{Failure Mode and Effect Analysis}
\label{TblFMEA}\\
\hline
\textbf{Component} & \textbf{Failure Modes} & \textbf{Effects of Failure} &
\textbf{Causes of Failure} & \textbf{Recommended Action} &
\textbf{SR / SRS Ref.} \\
\hline
\endfirsthead
\hline
\textbf{Component} & \textbf{Failure Modes} & \textbf{Effects of Failure} &
\textbf{Causes of Failure} & \textbf{Recommended Action} &
\textbf{SR / SRS Ref.} \\
\hline
\endhead
\hline
\endfoot

Intent Capture &
Extracted requirements are incomplete or incorrect &
Wrong or incomplete proposal; incorrect code if approved &
a. Ambiguous developer prompt\newline
b. LLM extraction error\newline
c. Schema mismatch &
a. Present proposal for mandatory human review before approval\newline
b. Allow edit/reject of individual items\newline
c. Validate extracted content against the intent schema &
a. SR1\newline
b. FR3--FR6\newline
H1-1 \\
\hline

Intent Capture &
Extraction fails or times out &
Developer cannot capture intent; workflow blocked &
a. Network/LLM outage\newline
b. Service rate limit\newline
c. Malformed response &
a. Notify the developer of extraction failure\newline
b. Allow retry from the last request\newline
c. Do not proceed to code generation while extraction is incomplete &
a. SR2\newline
b. FR7\newline
H1-2 \\
\hline

Review and Approval &
Proposal approved without meaningful review &
Incorrect intent enters the intent model; downstream code follows bad
requirements &
a. Developer rubber-stamps proposal\newline
b. UI makes approve too easy relative to review &
a. Require explicit approve/reject action per proposal\newline
b. Display structured requirements/constraints clearly before approve &
a. SR1\newline
b. FR4--FR5\newline
H2-1 \\
\hline

Review and Approval &
Developer cannot approve, edit, or reject &
Workflow stuck; temptation to bypass the tool and generate code directly &
a. UI/extension crash\newline
b. Corrupted draft proposal &
a. Persist drafts locally and recover on restart\newline
b. Allow discarding a draft and re-extracting &
a. SR3\newline
b. FR5--FR6\newline
H2-2 \\
\hline

Code Generation Control &
Code is generated before approval &
AI implements unapproved intent; loss of human control (G5) &
a. Gate not enforced by assistant integration\newline
b. Developer bypasses tool &
a. Block assistant implementation until approve/reject is recorded\newline
b. Document that bypassing the gate is unsupported and breaks the safety case &
a. SR4\newline
b. FR7--FR8\newline
H3-1 \\
\hline

Code Generation Control &
Approval decision not communicated to assistant / store &
Assistant implements wrong decision; intent model out of sync &
a. Integration message drop\newline
b. Partial write to store &
a. Treat decision recording as atomic with the gate signal\newline
b. Report failure if either side does not acknowledge &
a. SR5\newline
b. FR8\newline
H3-2 \\
\hline

Intent Model Store &
Intent model corrupted, truncated, or unreadable &
Loss of source of truth; verification and traceability fail &
a. Disk/write failure\newline
b. Concurrent write\newline
c. Invalid schema write &
a. Validate on write against schema\newline
b. Keep recoverable versions / backups of the intent model &
a. SR6\newline
b. FR9\newline
H4-1 \\
\hline

Intent Model Store &
Intent model modified without recorded approval &
Tampered lineage; false sense of reviewed requirements &
a. Manual file edit\newline
b. Bug writes approved state incorrectly &
a. Record provenance for approved changes\newline
b. Prefer append/versioned updates over silent overwrite &
a. SR7\newline
b. FR9, FR11\newline
H4-2 \\
\hline

Traceability &
Trace links missing, stale, or unjustified &
Cannot explain why code exists; maintainability/scalability loss &
a. Code changed outside tool\newline
b. Link update failure after refactor &
a. Detect and report broken or missing links on verification\newline
b. Require justification when creating links &
a. SR8\newline
b. FR9, FR11\newline
H5-1 \\
\hline

Change Verification &
Verification reports pass for a violating change &
Unsafe/incorrect AI change accepted (false confidence) &
a. Incomplete compiled checks\newline
b. Stale intent model used for check\newline
c. Heuristic/LLM checker error &
a. Compile approved constraints into executable checks where possible\newline
b. Always verify against the latest approved intent model\newline
c. Present verification evidence for developer review &
a. SR9\newline
b. FR10, FR13\newline
H6-1 \\
\hline

Change Verification &
Verification unavailable or fails closed poorly &
Developer either blocked indefinitely or skips checks &
a. Checker crash\newline
b. Missing tests/checks &
a. Fail visibly with an error status (do not silent-pass)\newline
b. Allow limited proceed only with explicit developer acknowledgement &
a. SR10\newline
b. FR10\newline
H6-2 \\
\hline

IDE Extension / UI &
Extension fails to load or loses workspace sync &
Developer cannot see proposals, status, or verification results &
a. IDE API change\newline
b. File-watcher failure &
a. Surface a clear degraded-mode banner\newline
b. Provide a command to reload/resync project state &
a. SR2\newline
H7-1 \\
\hline

External LLM / Network &
Project content leaked to external services &
Confidential code/requirements exposed; reputational/legal loss &
a. Full-file prompts to hosted LLM\newline
b. Logging of prompts/responses &
a. Minimize outbound context; prefer local processing where feasible\newline
b. Document what leaves the machine; allow opt-out of external calls &
a. SR11\newline
H8-1 \\
\hline

External LLM / Network &
Network unavailable during extraction/verification &
Features depending on remote LLM unavailable &
a. Offline work\newline
b. Firewall/proxy &
a. Notify user of offline/degraded mode\newline
b. Keep local intent model usable for review of already-captured intent &
a. SR2\newline
H8-2 \\
\hline

Version Control &
Intent model not versioned or diverges from repo history &
Cannot reconstruct past intent; rationale reconstruction fails &
a. Intent model not committed\newline
b. Merge conflict mishandled &
a. Store intent model in-repo and document commit expectations\newline
b. Detect uncommitted intent-model changes in the UI &
a. SR12\newline
b. FR9\newline
H9-1 \\
\hline

\end{longtable}
\end{small}
\end{landscape}

\section{Safety and Security Requirements}
\label{sec:SR}

The following requirements were identified through the FMEA. They should also be
added to the SRS (for example under Safety-Critical, Robustness, Security, or
Operational requirements). Existing functional requirements that already mitigate
a hazard are cited in Table~\ref{TblFMEA}.

\noindent\textbf{SR1.} The system shall require explicit developer review and an
approve, edit, or reject decision on every extracted proposal before the
proposal can update the intent model or authorize code generation.\\
\textit{Rationale:} Incorrect extraction must not silently become approved
intent.\\
\textit{Associated hazards:} H1-1, H2-1\\
\textit{Related SRS:} FR3--FR6, G1, G5

\medskip
\noindent\textbf{SR2.} The system shall notify the developer when intent capture,
verification, or IDE sync fails or is unavailable, including network/LLM
outages.\\
\textit{Rationale:} Silent failure encourages unsafe bypass of the tool.\\
\textit{Associated hazards:} H1-2, H7-1, H8-2

\medskip
\noindent\textbf{SR3.} The system shall persist draft proposals locally so that an
interrupted review can be resumed or safely discarded and re-run.\\
\textit{Rationale:} Lost drafts push developers to skip the review workflow.\\
\textit{Associated hazards:} H2-2\\
\textit{Related SRS:} FR5--FR6

\medskip
\noindent\textbf{SR4.} The system shall prevent the integrated AI coding assistant from
implementing a request until the corresponding proposal has an explicit
approve or reject decision recorded by the developer.\\
\textit{Rationale:} Human-in-the-loop control is a primary project goal
(G5).\\
\textit{Associated hazards:} H3-1\\
\textit{Related SRS:} FR7--FR8

\medskip
\noindent\textbf{SR5.} The system shall treat recording of the developer's decision in the
intent-management store and signalling that decision to the AI coding assistant
as a single successful transaction, or report failure if either side does not
complete.\\
\textit{Rationale:} Partial updates desynchronize the gate and the source of
truth.\\
\textit{Associated hazards:} H3-2\\
\textit{Related SRS:} FR8

\medskip
\noindent\textbf{SR6.} The system shall validate intent-model writes against the defined
schema and retain recoverable prior versions of the intent model.\\
\textit{Rationale:} Corruption or truncation of the intent model is a loss of
mission for the tool.\\
\textit{Associated hazards:} H4-1\\
\textit{Related SRS:} FR9

\medskip
\noindent\textbf{SR7.} The system shall record provenance for approved intent-model
changes (who/what approved, and from which proposal/request) and shall not
silently overwrite approved intent without a new recorded decision.\\
\textit{Rationale:} Untracked edits break requirement lineage.\\
\textit{Associated hazards:} H4-2\\
\textit{Related SRS:} FR9, FR11

\medskip
\noindent\textbf{SR8.} The system shall detect and report missing, broken, or unjustified
traceability links between intent-model elements and associated code or
tests.\\
\textit{Rationale:} Traceability is required for maintainability and for
explaining AI-generated structure.\\
\textit{Associated hazards:} H5-1\\
\textit{Related SRS:} FR9, FR11, G3

\medskip
\noindent\textbf{SR9.} The system shall verify AI-proposed or generated changes against the
latest approved intent model, prefer executable compiled checks where
available, and present verification evidence to the developer rather than a
bare pass/fail with no support.\\
\textit{Rationale:} False passes create over-trust in unsafe changes.\\
\textit{Associated hazards:} H6-1\\
\textit{Related SRS:} FR10, FR13, G4

\medskip
\noindent\textbf{SR10.} The system shall not treat a verification infrastructure failure as
a successful pass; verification errors shall be surfaced explicitly, and any
continue-anyway path shall require explicit developer acknowledgement.\\
\textit{Rationale:} Silent pass-on-error is an unacceptable hazard control
failure.\\
\textit{Associated hazards:} H6-2

\medskip
\noindent\textbf{SR11.} The system shall minimize project content sent to external LLM or
logging services, document what data may leave the machine, and provide a way
to disable or avoid external calls when the developer requires local-only
operation.\\
\textit{Rationale:} Leakage of private code or requirements is a stakeholder
loss.\\
\textit{Associated hazards:} H8-1

\medskip
\noindent\textbf{SR12.} The system shall keep the intent model as a versionable project
artifact compatible with Git and shall warn when intent-model changes are
uncommitted relative to the workspace.\\
\textit{Rationale:} Without versioning, past intent and rationale cannot be
reconstructed.\\
\textit{Associated hazards:} H9-1\\
\textit{Related SRS:} FR9

\section{Roadmap}

Not all safety and security requirements can be fully implemented within the
capstone timeline. Priority for Revision~0 / Proof of Concept follows the
project's stated primary risk: constructing and maintaining a trustworthy
mapping between intent and implementation under human approval.

\begin{itemize}
  \item \textbf{In scope for the capstone timeline (Rev~0 / PoC focus):}
  SR1, SR4, SR5, SR6, SR9, and the related SRS requirements FR3--FR10 that
  already encode capture, approval, gating, intent persistence, and
  verification.
  \item \textbf{Targeted for later / stretch within the course:}
  SR3, SR7, SR8, SR10, SR12 (stronger provenance, trace-break detection,
  continue-anyway acknowledgement UX, and Git hygiene warnings).
  \item \textbf{Likely future work beyond the course if not reached:}
  SR2 polish across all failure surfaces, and SR11 (local-only / privacy
  controls for external LLM traffic).
\end{itemize}

\newpage{}

\section{Appendix --- Generative AI Use Disclosure}

\subsection{AI Tools Used}

\begin{itemize}
  \item Claude (Anthropic), Claude Sonnet 4.6
  \item Cursor Agent (Composer), used while drafting this document
\end{itemize}

\subsection{How and Where Used}

\begin{itemize}
  \item Drafting and structuring Sections~1--7 of this hazard analysis from the
  project SRS and a prior-year example (Sayyara), including the FMEA table and
  SR1--SR12 wording.
  \item Spelling, grammar, and clarity edits on drafted prose.
\end{itemize}

\subsection{Verification of AI-Generated Content}

All AI-assisted text was treated as a draft. Team members will review each
FMEA row and safety requirement against the actual product goals and SRS before
submission. Hazard identifiers, requirement IDs, and claims about intended
behaviour were checked against \texttt{docs/SRS-Volere/SRS.tex}. The team remains
responsible for the accuracy and completeness of this document.

\newpage{}

\section{Appendix --- Reflection}

\wss{Not required for CAS 741}

\input{../Reflection.text}

\begin{enumerate}
    \item What went well while writing this deliverable?
    \item What pain points did you experience during this deliverable, and how
    did you resolve them?
    \item Which of your listed risks had your team thought of before this
    deliverable, and which did you think of while doing this deliverable? For
    the latter ones (ones you thought of while doing the Hazard Analysis), how
    did they come about?
    \item Other than the risk of physical harm (some projects may not have any
    appreciable risks of this form), list at least 2 other types of risk in
    software products. Why are they important to consider?
\end{enumerate}

\begin{thebibliography}{9}
\bibitem{STPA2018}
Nancy G. Leveson and John P. Thomas.
\newblock \emph{STPA Handbook}.
\newblock March 2018.
\end{thebibliography}

\end{document}
